%% ============================================================================
%%  SODE-Guard Revision — Code Changelog (v3 -> v4)
%%
%%  Companion to the IEEE Access resubmission (Access-2026-27603).
%%  Compile with pdflatex (self-contained).
%% ============================================================================
\documentclass[11pt,letterpaper]{article}

\usepackage[margin=1in]{geometry}
\usepackage[T1]{fontenc}
\usepackage{lmodern}
\usepackage{amsmath,amssymb}
\usepackage{booktabs}
\usepackage{longtable}
\usepackage{tabularx}
\usepackage{enumitem}
\usepackage{xcolor}
\usepackage[colorlinks=true,linkcolor=blue!60!black,urlcolor=blue!60!black]{hyperref}
\usepackage{fancyhdr}
\usepackage{titlesec}

\setlength{\parindent}{0pt}
\setlength{\parskip}{6pt}
\pagestyle{fancy}
\fancyhf{}
\lhead{\small SODE-Guard v4 — Code changelog}
\rhead{\small \texttt{rogerpanel/CV/SODE-Guard}}
\cfoot{\thepage}
\titleformat{\section}{\normalfont\Large\bfseries\color{blue!60!black}}{\thesection}{0.7em}{}

\title{\vspace{-1.2cm}\textbf{SODE-Guard Revision — Code Changelog}\\[4pt]
       \normalsize v3 $\rightarrow$ v4 (IEEE Access MS \texttt{Access-2026-27603})}
\author{R.\ N.\ Anaedevha, A.\ G.\ Trofimov, and Y.\ V.\ Borodachev\\
        \normalsize National Research Nuclear University MEPhI, Moscow, Russia}
\date{\normalsize \today}

\begin{document}
\maketitle
\thispagestyle{fancy}
\vspace{-0.6cm}

\section*{Overview}
This document enumerates every code addition, modification, and new
artefact introduced by the v3 $\rightarrow$ v4 revision. All paths are
relative to the reproducibility root at
\url{https://github.com/rogerpanel/CV/tree/main/SODE-Guard}.

%% ============================================================================
\section{New source modules}
%% ============================================================================
\begin{longtable}{@{}p{0.42\linewidth}p{0.16\linewidth}p{0.36\linewidth}@{}}
\toprule
\textbf{Module} & \textbf{Concern} & \textbf{Purpose} \\
\midrule
\endhead
\texttt{src/theory/lipschitz.py} & R1.1, R2.2 & Empirical Hoeffding $L_g$ certification \\
\texttt{src/theory/chaos\_degree.py} & R2.4 & Data-driven Hermite chaos-degree estimator \\
\texttt{src/theory/carbery\_wright.py} & R2.2, R2.3 & Corrected margin-stability + decision-flip bounds \\
\texttt{src/theory/pac\_bayes.py} & R1.6 & Maurer PAC-Bayes-kl certificate \\
\texttt{src/theory/pseudoinverse.py} & R1.5, R2.9 & Moore--Penrose pseudo-inverse for $128\times 16$ diffusion \\
\texttt{src/theory/bel\_estimator.py} & R2.8 & BEL vs finite-difference diagnostic sweep \\
\texttt{src/attacks/semantic/feasibility.py} & R2.11 & Feature-level feasibility projector \\
\texttt{src/attacks/semantic/constrained\_attacks.py} & R2.11 & Problem-space PGD / C\&W \\
\texttt{src/data/splits\_extra/temporal.py} & R2.13 & Time-ordered split \\
\texttt{src/data/splits\_extra/host\_disjoint.py} & R2.13 & Host-disjoint split \\
\texttt{src/data/splits\_extra/scenario\_disjoint.py} & R2.13 & Scenario-disjoint split \\
\texttt{src/data/splits\_extra/dedup.py} & R2.12 & Deduplication + leakage report \\
\texttt{src/evaluation/reliability/smoothing\_sweep.py} & R1.3 & Randomised-smoothing $\sigma$ sensitivity \\
\texttt{src/evaluation/reliability/commercial\_context.py} & R1.2 & Baseline threat-model taxonomy \\
\bottomrule
\end{longtable}

%% ============================================================================
\section{Modified modules}
%% ============================================================================
\begin{itemize}[leftmargin=1.4em]
  \item \texttt{src/models/sode\_guard.py}: \texttt{forward\_with\_paths} now
        also exposes the Brownian summary needed by
        \texttt{theory.chaos\_degree}.
  \item \texttt{src/regularizers/anti\_concentration.py}: docstring
        corrected to reflect the direction-fixed Carbery--Wright bound;
        code unchanged because it uses the empirical $\|m\|_2$
        denominator.
  \item \texttt{src/evaluation/certificate.py}: consumes
        \texttt{theory.lipschitz.LipschitzCertificate} instead of a
        hand-picked Lipschitz probe.
\end{itemize}

%% ============================================================================
\section{New documentation and manuscript files}
%% ============================================================================
\begin{itemize}[leftmargin=1.4em]
  \item \texttt{docs/response\_to\_reviewers/RESPONSE.md} (Markdown)
        and its LaTeX equivalent
        \texttt{manuscript/Response\_to\_Reviewers.tex}.
  \item \texttt{docs/response\_to\_reviewers/CHANGELOG.md} (Markdown)
        and its LaTeX equivalent
        \texttt{manuscript/Code\_Changelog.tex} (this file).
  \item \texttt{docs/packet\_semantics.md} (Markdown) and
        \texttt{manuscript/Packet\_Semantics.tex}.
  \item \texttt{manuscript/SODEGuard\_v4\_Manuscript.tex} (revised
        submission LaTeX).
  \item \texttt{manuscript/highlights\_v3\_to\_v4.md} and
        \texttt{manuscript/Highlights\_v3\_to\_v4.tex} (change list
        mapped to line ranges for the Highlighted-PDF submission).
\end{itemize}

%% ============================================================================
\section{New tests}
%% ============================================================================
\begin{itemize}[leftmargin=1.4em]
  \item \texttt{tests/test\_lipschitz.py}
  \item \texttt{tests/test\_chaos\_degree.py}
  \item \texttt{tests/test\_carbery\_wright.py}
  \item \texttt{tests/test\_pac\_bayes.py}
  \item \texttt{tests/test\_pseudoinverse.py}
  \item \texttt{tests/test\_feasibility\_projector.py}
  \item \texttt{tests/test\_split\_protocols.py}
  \item \texttt{tests/test\_dedup\_leakage.py}
\end{itemize}

%% ============================================================================
\section{New scripts}
%% ============================================================================
\begin{itemize}[leftmargin=1.4em]
  \item \texttt{scripts/reproduce\_revision.sh} — re-runs Tables 5--11
        of the revised manuscript on the local hardware.
  \item \texttt{scripts/regenerate\_figures.py} — regenerates
        Figures~4b ($\sigma$ sweep) and 6 (commercial-baseline
        reference plot).
\end{itemize}

\end{document}
